Después de completar todo el desarrollo de Thary, se identificaron diferentes oportunidades de mejora que podrían hacer que el sistema mejore su desempeño, seguridad, eficiencia y en general, ser un mejor producto para el usuario final. Aunque dichas oportunidades fueron identificadas, no pudieron ser implementadas en esta versión del sistema debido al alcance del proyecto y demas limitaciones. Aun así, se consideran aportes valiosos que valen la pena ser documentados para futuras actualizaciones y versiones del sistema.


\subsubsection{Realización de validación con usuarios}

La validación de usabilidad fue realizada y documentada en el capítulo 4, sin embargo, esta consto con un carácter informal y con un número bastante reducido de participantes. Es por eso que se propone como trabajo futuro realizar un estudio de usabilidad más formal, el cual cuente con un mayor número de participantes con el perfil del usuario objetivo (es decir, personas relacionadas con la administración de productos en inventarios, ya sea con el rol de personal de inventario o de administrador). También convendría incluir métricas cuantitativas de satisfacción y facilidad de uso como implementar cuestionarios estandarizados como el System Usability Scale, más allá de solo pedir que prueben el sistema y recopilar feedback.


\subsubsection{Conservación del historial de artefactos por corrida}

El sistema está configurado de manera que cada vez que se procesa un archivo, MySQL guarda un registro nuevo sin borrar los anteriores, sin embargo, los archivos asociados a esa corrida como el modelo entrenado y las gráficas, se guardan siempre en la misma carpeta, por lo que se sobrescriben con cada nueva ejecución. Esto significa que solo se guardan los archivos de la ejecución más reciente, aunque en la base de datos si exista el historial completo de todas las ejecuciones anteriores.
Es por esto que seria valioso encontrar alguna manera para poder guardar diferentes archivos por cada corrida, para en el futuro, poder comparar como se comportaba el modelo entrenado en distintos momentos, y no solo ver sus métricas ya calculadas.

\subsubsection{Validación en otros dominios de inventario}

Uno de los requerimientos no funcionales más importantes de Thary es que el sistema debe de poder ser utilizado para adaptarse a distintos dominios de inventarios. Sin embargo, a la hora de probar y validar el sistema, se utilizó como caso principal un único dataset real, correspondiente al dominio de la farmacia hospitalaria y al consumo de medicamentos e insumos médicos de la Clínica Nuestra Señora de los Remedios durante el periodo comprendido entre julio de 2021 y junio de 2022. Este dataset fue utilizado para probar las diferentes funcionalidades del sistema, debido a que es el conjunto de datos real más completo con el que se contaba y el único que representaba de manera realista el comportamiento de la demanda de un inventario.

Adicionalmente, se realizaron pruebas exploratorias con otros datasets públicos obtenidos de Kaggle, pertenecientes a diferentes contextos de inventario. Estos datasets permitieron observar el comportamiento del sistema en otros escenarios y complementar la evaluación realizada con el caso principal. Sin embargo, al tratarse de datos públicos y no de datos reales obtenidos directamente de organizaciones de dichos dominios, estos resultados no permiten confirmar completamente la adaptabilidad del sistema a otros contextos reales.

Por esta razón, como trabajo futuro se plantea realizar pruebas con datasets reales de otros dominios de inventario, como retail, ferretería, alimentos, entre otros. Esto permitiría evaluar de manera más completa la capacidad de adaptación de Thary y comprobar si su funcionamiento y desempeño se mantienen cuando se utilizan datos reales con características diferentes a las del dominio de la farmacia hospitalaria.

\subsubsection{Exploración de modelos de aprendizaje profundo}

Durante la revisión de literatura se exploraron diferentes técnicas de Deep Learning con las que es posible realizar predicciones de series de tiempo. Entre ellas, se descartaron LSTM y GRU como modelos candidatos para el sistema Thary, aun así, en un trabajo futuro se podría rescatar la posibilidad de integrar alguna de estas técnicas como una cuarta o quinta alternativa al mecanismo de selección de mejor modelo por producto. Principalmente en escenarios donde se cuente con un historial de consumo más extenso al utilizado en el caso de estudio actual, dado que este tipo de modelos suele requerir mayor cantidad de datos para superar a técnicas más simples.

\subsubsection{Reevaluación del asistente conversacional}

Durante una primera versión de Thary, se implementó un asistente conversacional basado en un modelo de lenguaje ejecutado localmente (Ollama), el cual fue posteriormente descartado como una de las funcionalidades del sistema en la versión final. Se podría considerar en el futuro reevaluar la reactivacion de esta funcionalidad, incluyendo un análisis más detallado de su valor real para el usuario final.

\subsubsection{Combinación de pronósticos con métodos más sofisticados}

La combinación de pronósticos implementada en Thary utiliza un único peso fijo, calculado por mínimos cuadrados sobre todo el periodo de backtest, que es el método más simple dentro de la literatura de combinación de pronósticos \cite{wang2023combinations}. Se propone como trabajo futuro explorar métodos más sofisticados, como pesos que varíen en el tiempo o combinaciones que incluyan los tres modelos candidatos en lugar de solo dos.

\subsection{Reentrenamiento de los modelos antes del pronóstico futuro}

Actualmente, los modelos de Regresión Lineal y XGBoost se entrenan una sola vez con los meses de entrenamiento, y esos mismos modelos se reutilizan tanto para elegir el campeón como para generar el pronóstico futuro. Esto quiere decir que, en el caso de estudio de la clínica, el pronóstico futuro se genera con modelos que solo aprendieron de octubre a diciembre de 2021, sin aprovechar los seis meses más recientes del historial. En este caso el impacto es pequeño, ya que la gran mayoría de los productos usan Media Móvil, que siempre toma los últimos meses reales, pero en catálogos como el de Store Item Demand Forecasting, donde XGBoost es el campeón en más de la mitad de las series, sí puede afectar la calidad del pronóstico. Se propone que, una vez elegido el campeón de cada producto, los modelos se vuelvan a entrenar con todos los meses disponibles antes de generar el pronóstico futuro.